\subsection{The leading quotient and the mixed numerical factor}
\label{subsec:selmer-mixed-count}

The numerical factor in the lower comparison has an elementary origin.
Retain the graded algebra, rational chart and three modules
$N_0,N_1,N_2$ above. Assume $A_{-2}$ and $A_2$ are finite over $B$, that
the chart preserves weights, has the augmentation as its constant term,
and is surjective on linear terms over $F$. Assume also that
$\varphi(F_*)$ has zero linear part. The chart variables have weights
$-2,-1,1,2$.
For a monomial, write $\nu_-$ and $\nu_+$ for the sums of its absolute
negative and positive weights, counted with exponents. Put
\[
 a(\boldsymbol d)=2d_0+4d_2,\qquad b(\boldsymbol d)=2d_1+2d_2.
\]

\begin{lemma}[The finite free leading piece]
\label{lem:selmer-leading-piece}
In degree $\boldsymbol d\ge0$, retaining only the terms of negative
weight $a(\boldsymbol d)$ in the chart defines a graded algebra map
\[
 \mathcal R\longrightarrow
 \bigoplus_{\boldsymbol d\ge0}
 F[z]_{a(\boldsymbol d),b(\boldsymbol d)}t^{\boldsymbol d},
\]
where the subscript specifies the pair $(\nu_-,\nu_+)$. Let
$C=\mathcal R/\mathfrak q$ be its image. Each $C_{\boldsymbol d}$ is
finite free over $\mathcal O$, with rank
$\dim_F F[z]_{a(\boldsymbol d),b(\boldsymbol d)}$.
The action of $B$ on $C$ is through the augmentation. The ideal
$\mathfrak q$ contains $J_\varphi$ and all degree-zero elements of
augmentation zero, and is saturated under nonzero scalars of
$\mathcal O$. If $f,g\in A_{-1}$ have the same chart linear part, then
$(f^2-g^2)t_0\in\mathfrak q$.
\end{lemma}

\begin{proof}
The minimum negative weights in $N_0,N_1,N_2$ are respectively at least
$2,0,4$. Indeed, a weight $-1$ series without a linear term has negative
weight at least two. These bounds add under products, proving that the
leading projection is multiplicative. A weight-zero series of
augmentation zero has negative weight at least one, so $B$ acts through
the augmentation. Thus each leading piece is finite over $\mathcal O$.
It lies in an $F$-vector space, and hence is torsion free and free.

To see full generic rank, choose rational homogeneous preimages of the
linear variables. Products of the weight $-1$ preimages supply every
quadratic monomial of pure negative weight two. A preimage of a weight
$-2$ variable has that variable as its linear part; subtracting those
quadratic products isolates it in leading weight two. Consequently
$A_{-2}\otimes F$ fills $F[z]_{2,0}$. The positive variables give
$A_2\otimes F$ filling $F[z]_{0,2}$. Every monomial of bidegree $(2u,2v)$
factors into $u$ negative weight-two blocks and $v$ positive weight-two
blocks, by pairing weight-one variables. Since
\[
 A_{-2}^{d_0+2d_2}A_2^{d_1+d_2}
 \subset N_0^{d_0}N_1^{d_1}N_2^{d_2},
\]
the leading image has full rank. The kernel properties follow from the
map into an $F$-algebra. Finally, $(f-g)(f+g)$ has negative weight at
least three, so its degree-$e_0$ leading projection is zero.
\end{proof}

Let the variable counts in weights $-1,-2,1,2$ be
$r,s,r'+1,s'+1$, respectively, with $r\ge1$ and $s,r',s'\ge0$.
Write $H(a,b)$ for the number of monomials of bidegree $(a,b)$, and set
it to zero if either index is negative. Define the backward difference
$\Delta_{(u,v)}H(a,b)=H(a,b)-H(a-u,b-v)$.

\begin{lemma}[Exact mixed count]
\label{lem:selmer-mixed-count}
For $a\ge r+3$ and $b\ge r'$, one has
\[
 \Delta_{(2,0)}^{r-1+s}\Delta_{(0,2)}^{r'+s'}
 \Delta_{(4,2)}H(a,b)=2^{r-1+r'}.
\]
\end{lemma}

\begin{proof}
The generating series of $H$ is
\[
 \frac1{(1-X)^r(1-X^2)^s(1-Y)^{r'+1}(1-Y^2)^{s'+1}}.
\]
Applying the backward differences multiplies it by the corresponding
factors $1-X^uY^v$, giving
\[
 \frac{(1+X)^{r-1}(1+Y)^{r'}(1-X^4Y^2)}
      {(1-X)(1-Y)(1-Y^2)}.
\]
Use $1-X^4Y^2=(1-Y^2)+Y^2(1-X^4)$. The first resulting summand is
$(1+X)^{r-1}(1+Y)^{r'}/((1-X)(1-Y))$, whose coefficient is
$2^{r-1+r'}$ in the stated range. The second has numerator
\[
 Y^2(1+X)^{r-1}(1+X+X^2+X^3)(1+Y)^{r'}
\]
and denominator $(1-Y)(1-Y^2)$. Its degree in $X$ is at most $r+2$,
so it contributes zero. This also covers $r=1$ and $r'=0$.
\end{proof}

For the Hilbert function
$h(\boldsymbol d)=\ell_{\mathcal O}(C_{\boldsymbol d}/\pi C_{\boldsymbol d})$,
Lemma~\ref{lem:selmer-leading-piece} gives
$h(\boldsymbol d)=H(2d_0+4d_2,2d_1+2d_2)$ for $\boldsymbol d\ge0$.
The length over $B$ is the same, since $B$ acts through the surjection
$x:B\to\mathcal O$.
The shifts by $e_0,e_1,e_2$ become $(2,0),(0,2),(4,2)$. Consequently
\[
 \Delta_0^{r-1+s}\Delta_1^{r'+s'}\Delta_2h(\boldsymbol d)
 =2^{r-1+r'}
\]
throughout the explicit box
\[
 d_0\ge r-1+s,\qquad d_1\ge r'+s',\qquad d_2\ge1.
\]
Every shifted multidegree then stays nonnegative. Extracting the factors
$\pi^{2c_i}$ in the remaining intersection argument gives the scalar
\[
 \left(\sum_i2c_i\right)2^{r-1+r'}=2^{r+r'}\sum_i c_i.
\]
The leading quotient and this scalar are now proved directly. The next
argument supplies positivity, quotient comparison and exact extraction
of the scalar factors, under the explicit admissibility condition
needed for the lower bound in \eqref{eq:selmer-euler-retained-inputs}.
